\chapter*{About this manual}
\addcontentsline{toc}{chapter}{About this manual}

This manual is for the people who use WOP in the workshop: the ones who have a part to turn and want a G-code program for the lathe. It explains what each menu, window and button does and how to get a program from the screen to the machine. It does not explain how WOP is built inside; that is in the documents for developers.

\section*{How to read it}
\begin{itemize}
  \item \textbf{Part I} tells you what WOP is, how to start it, and how the main window is laid out.
  \item \textbf{Part II} goes through the menus, one chapter for each menu, in the order of the menu bar.
  \item \textbf{Part III} follows the work: the project, the blocks, contours, tools, workholding, machines, checking and export.
  \item \textbf{Part IV} is for looking things up: keyboard shortcuts, files and folders, a glossary, and what to do when something goes wrong.
\end{itemize}

\section*{How things are written}
\begin{center}
\begin{tabularx}{0.9\textwidth}{@{}l Y@{}}
\toprule
\textbf{Looks like} & \textbf{Means} \\
\midrule
\menupath{File}{Open...} & A menu and an entry in it. \\
\button{Close} & A button in a window. \\
\field{Outer diameter} & The name of a field, a list or a check box. \\
\key{Ctrl}+\key{S} & Keys on the keyboard, pressed together. \\
\file{project.wop.json} & A file name, or text you type. \\
\bottomrule
\end{tabularx}
\end{center}

\begin{notebox}
A note adds something that is good to know.
\end{notebox}
\begin{tipbox}
A tip makes the work easier.
\end{tipbox}
\begin{warningbox}
A warning is about something that can cost you work, a tool or a part.
\end{warningbox}

\section*{Status of this draft}
This manual is machine generated.

WOP is still being developed, and so is this manual. Every chapter is written, but it was written from the program and has not yet been checked by many readers. Where something is not confirmed on a real machine, the text says so. If you find a mistake or something that is missing, please tell the developer. Pictures are taken by a tool from the program itself, so they follow the program when it changes.
